\lstset{style=Python}
\chapter{Anschluss von Datenbank an Python-Game}
\section{Einführung}

Damit ein Game auch einen Highscore speichern kann, welcher beim nächsten Spielstart angezeigt wird, muss es an eine Datenbank angeschlossen werden, in der die Daten permanent gespeichert werden können. Dies kann in Python mithilfe der Library \lstinline|sqlite3| geschehen, welche die Interaktion mit einer SQLite-Datenbank ermöglicht. SQLite ist eine leichtgewichtige, serverlose Datenbank, welche sich hervorragend für lokale Anwendungen wie unser Python-Game eignet.

In diesem Kapitel werden wir unser Python-Game an eine Datenbank anschliessen, um neue Funktionalitäten wie ein Login-Fenster, einen Highscore und weitere Dinge zu ermöglichen. Hierzu möchten wir als Erstes folgende Tabellen erstellen:


\begin{figure}[H]
    \centering
    \begin{tikzpicture}[
            table/.style={
                    matrix of nodes,
                    nodes={anchor=west, minimum height=1em, text depth=0.2em, text height=1em, align=left},
                    column 1/.style={right delimiter=|},
                    row sep=0pt,
                    column sep=2pt
                }
        ]

        % Tabelle scores
        \matrix (scores) [table] {
                                   & \textbf{Spalte}      & \textbf{Typ}      \\\midrule\midrule
            \emoji{key}            & \texttt{id}          & \texttt{INTEGER}  \\\midrule
            \emoji{up-right-arrow} & \texttt{user\_id}    & \texttt{INTEGER}  \\\midrule
                                   & \texttt{score}       & \texttt{INTEGER}  \\\midrule
                                   & \texttt{date\_score} & \texttt{DATETIME} \\
        };

        % Tabelle users
        \matrix (users) [table, right=5cm of scores] {
                        & \textbf{Spalte}   & \textbf{Typ}         \\\midrule\midrule
            \emoji{key} & \texttt{id}       & \texttt{INTEGER}     \\\midrule
                        & \texttt{user}     & \texttt{VARCHAR(40)} \\\midrule
                        & \texttt{password} & \texttt{VARCHAR(40)} \\
        };

        \node(userstitle)[above=5.5em of users.center, font=\bfseries, anchor=center]{Tabelle \texttt{users}};
        \node(scorestitle)[above=7em of scores.center, font=\bfseries, anchor=center]{Tabelle \texttt{scores}};

        % Pfeil von user_id (Zeile 3 in scores) zu id (Zeile 2 in users)
        \draw [->,>=latex, thick] ([yshift=10pt]scores-3-3.east) -- +(40pt,0) -| ([xshift=-40pt,yshift=10pt]users-2-1.west) -- ([yshift=10pt]users-2-1.west);

    \end{tikzpicture}
\end{figure}

Die Tabelle \texttt{users} speichert alle Benutzer-Daten für das Spiel, d.h., die Benutzer, welche sich in das Spiel eingeloggt haben. Die Tabelle \texttt{scores} speichert die Spielresultate jedes Spiels.

\begin{myexercise}[\texttt{users}-Tabelle erstellen]

    Erstellen Sie ein Programm in Python, um die Tabelle \texttt{users} in der Datenbank \texttt{game.db} zu erstellen. Die Datenbank wird mit dem unten stehenden Code automatisch angelegt, sollte sie nicht bereits existieren.

    \begin{lstlisting}
from sqlite3 import *
from prettytable import *   # Befehle zur Ausgabe von Tabellen

# Datenbank game.db (wenn nötig) erzeugen und verbinden
with connect("game.db") as con:
    # Datenbank-Cursor (=Lese/Schreibestift) holen
    cursor = con.cursor()

    #TODO Mehrzeiliger SQL-Code zur Erstellung einer Tabelle
    sql = """
        CREATE TABLE IF NOT EXISTS users (
           ...
    """

    # SQL ausführen
    cursor.execute(sql)

    print("New table X was created!")

\end{lstlisting}

    \textbf{Tipp}: Primärschlüssel müssen erstellt werden mit dem Befehl \lstinline[style=sql]|PRIMARY KEY AUTOINCREMENT|:
    \begin{lstlisting}[style=sql]
CREATE TABLE IF NOT EXISTS users(
	id INTEGER PRIMARY KEY AUTOINCREMENT, ...
)
\end{lstlisting}
    Mit dem Zusatz \lstinline[style=sql]|AUTOINCREMENT| werden die Werte der Spalte \texttt{id} automatisch bei jedem Eintrag um 1 vergrössert, so dass kein Wert für \texttt{id} beim \lstinline[style=sql]|INSERT|-Befehl mitgegeben werden muss.
    \begin{myanswer}
        \lstinputlisting[caption={SQLite-Tabelle erstellen}]{Grundlagen_Info/06_Datenbanken/Skript/Code/sqlite_create_table.py}
    \end{myanswer}
\end{myexercise}



\begin{myexample}[SQL-Tabellen oder -Werte in Python ausgeben]
    Resultate von SQL-Abfragen können auf zwei Arten ausgegeben werden:
    \begin{enumerate}
        \item Einzelne Werte ausgeben:
              \begin{lstlisting}
cursor.execute("SELECT * FROM scores")
resultset = cursor.fetchall()
# gib die dritte Spalte (2) der ersten Zeile (0) aus
print(resultset[0][2])
\end{lstlisting}

        \item Um eine ganze Tabelle, schön dargestellt auszugeben, muss die Library \lstinline|prettytable| geladen und wie folgt verwendet werden:
              \begin{lstlisting}
from prettytable import *

# ...

cursor.execute("SELECT * FROM scores")
printTable(cursor)
\end{lstlisting}
    \end{enumerate}
\end{myexample}

\begin{myexercise}[Benutzer-Registrierung]
    Schreiben Sie ein Programm, um Sie mittels \lstinline|input("...")| nach einem Benutzernamen und Passwort zu fragen. Beides sollte danach mit \lstinline[style=sql]|INSERT| in der Tabelle \texttt{users} abgespeichert werden. Überprüfen Sie am Ende das Resultat mit folgendem Code:
    \begin{lstlisting}
sql = "SELECT * FROM users"
cursor.execute(sql)
printTable(cursor)
\end{lstlisting}
    \textbf{Tipp}: Mit dem Text-Befehl \lstinline|.format()| lassen sich die Variablen-Werte innerhalb der Platzhalter \lstinline|{}| in einen Text einfügen.

    \begin{lstlisting}
sql = """INSERT INTO users (user, password)
        VALUES('{}', '{}')""".format(username, password)
\end{lstlisting}
    Die drei Anführungszeichen ermöglichen es, Texte über mehrere Zeilen zu schreiben.
    \begin{myanswer}
        \begin{lstlisting}
from sqlite3 import *
from prettytable import * # Befehle zur Ausgabe von Tabellen

# Datenbank game.db (wenn nötig) erzeugen und verbinden
with connect("game.db") as con:
    # Datenbank-Cursor (=Lese/Schreibestift) holen
    cursor = con.cursor()

    # Benutzername und Passwort abfragen
    username = input("Enter username")
    password = input("Enter password")

    # SQL-Code zum Einfügen eines neuen Datensatzes
    sql = """
        INSERT INTO users (user, password)
        VALUES
            ('{}', '{}')
    """.format(username, password)

    # SQL ausführen
    cursor.execute(sql)

    # Tabelle ausgeben
    sql = "SELECT * FROM users"
    cursor.execute(sql)
    printTable(cursor)

\end{lstlisting}
    \end{myanswer}
\end{myexercise}

\begin{mychallenge}[SQL-Injection verhindern]
    Das obige Programm fügt die Benutzereingaben mit \lstinline|.format()| direkt in den SQL-Befehl ein. Es ist dadurch anfällig für sogenannte SQL-Injections: Böswillige BenutzerInnen können ihre Eingabe so wählen, dass sich der SQL-Befehl selbst verändert. Besonders gefährlich ist dies bei einem Login, der Benutzername und Passwort in einer einzigen Abfrage prüft:
\begin{lstlisting}
sql = """
    SELECT * FROM users
    WHERE user = '{}' AND password = '{}'
""".format(username, password)
\end{lstlisting}
    Gibt jemand einen bestehenden Benutzernamen, z.B. \texttt{alice}, und als Passwort den Text \lstinline|' OR '1'='1| ein, führt die Datenbank folgenden Befehl aus:
\begin{lstlisting}[style=sql]
SELECT * FROM users
WHERE user = 'alice' AND password = '' OR '1'='1'
\end{lstlisting}
    Weil \lstinline[style=sql]|AND| vor \lstinline[style=sql]|OR| ausgewertet wird und \lstinline[style=sql]|'1'='1'| immer wahr ist, liefert die Abfrage alle Benutzer zurück. Ein Programm, das bei einem Treffer den Login erlaubt, lässt die Person also ohne Passwort herein. Angriffe mit einem zweiten Befehl wie \lstinline[style=sql]|'); DROP TABLE users; --| scheitern in Python an \lstinline|cursor.execute()|, weil diese Methode nur einen einzigen SQL-Befehl ausführt. Mit \lstinline|cursor.executescript()| oder in anderen Programmiersprachen funktionieren auch sie.

    Verhindern Sie SQL-Injections, indem Sie die Benutzereingaben mit einem sogenannten \textit{Prepared Statement} in den SQL-Befehl einfügen:
    \begin{lstlisting}
sql = """
    INSERT INTO users (user, password)
    VALUES
        (?, ?)
"""
cursor.execute(sql, (username, password))
\end{lstlisting}
Die Fragezeichen \lstinline|?| dienen als Platzhalter für die Werte, welche in der zweiten Zeile als Tupel übergeben werden. Auf diese Weise wird sichergestellt, dass die Benutzereingaben nicht als SQL-Befehle interpretiert werden können, sondern nur als Werte in die Tabelle eingefügt werden.
    \begin{myanswer}
        \begin{lstlisting}
from sqlite3 import *
from prettytable import * # Befehle zur Ausgabe von Tabellen

# Datenbank game.db (wenn nötig) erzeugen und verbinden
with connect("game.db") as con:
    # Datenbank-Cursor (=Lese/Schreibestift) holen
    cursor = con.cursor()

    # Benutzername und Passwort abfragen
    username = input("Enter username")
    password = input("Enter password")

    # SQL-Code zum Einfügen eines neuen Datensatzes
    sql = """
        INSERT INTO users (user, password)
        VALUES
            (?, ?)
    """
    cursor.execute(sql, (username, password))

    # Tabelle ausgeben
    sql = "SELECT * FROM users"
    cursor.execute(sql)
    printTable(cursor)
\end{lstlisting}
    \end{myanswer}

\end{mychallenge}

\begin{mychallenge}[Sicherheits-Überprüfung]
    Fügen Sie eine Bedingung ein, wonach eine Aufforderung erscheint, einen anderen Benutzernamen einzugeben, sollte es einen bestimmten Benutzer bereits geben.
    \begin{myanswer}
        \begin{lstlisting}
from sqlite3 import *

with connect("game.db") as con:
    # Datenbank-Cursor (=Lese/Schreibestift) holen
    cursor = con.cursor()

    # Benutzername abfragen
    while True:
        username = input("Enter username")
        sql = """
            SELECT *
            FROM users
            WHERE user = ?
        """
        cursor.execute(sql, (username,))
        resultset = cursor.fetchall()
        if len(resultset) == 0:
            # Noch unbenutzter Benutzername --> alles OK
            break
        else:
            print("This username is already taken. Choose a different one!")

    password = input("Enter password")

    # Restlicher Code

\end{lstlisting}
    \end{myanswer}
\end{mychallenge}

\begin{mychallenge}[Passwort hashen]
    In der Praxis werden Passwörter nicht im Klartext in Datenbanken gespeichert, sondern als \textit{Hashwert}. Eine Hashfunktion ist eine kryptografische Einwegfunktion: Sie berechnet aus dem Passwort einen Wert fester Länge, aus dem sich das Passwort nicht zurückrechnen lässt. Im Unterschied zu einer Verschlüsselung gibt es keinen Schlüssel und keine Entschlüsselung. Beim Login wird das eingegebene Passwort erneut gehasht und mit dem gespeicherten Hashwert verglichen. So kann nicht einmal der Administrator einer Webseite, welcher Zugriff auf die Datenbank hat, die Passwörter der BenutzerInnen lesen. Sie können anstelle des Passworts die Ausgabe der SHA-1-Hashfunktion des Passworts wie folgt speichern:
    \begin{lstlisting}
import hashlib
password = 'helloworld!'              # Ihr Klartext-Passwort zum "hashen"
hash_object = hashlib.sha1()          # SHA-1-Hash-Objekt erstellen
hash_object.update(password.encode()) # Hash-Objekt updaten mit der Byte-Darstellung des Passworts
hex_dig = hash_object.hexdigest()     # Hexadezimale Darstellung erhalten
print("SHA-1 Hash:", hex_dig)         # Hash ausgeben
\end{lstlisting}
    \textbf{Hinweis}: SHA-1 eignet sich für diese Übung, ist für echte Passwörter aber zu schnell berechnet. Angreifer können damit Milliarden von Passwörtern pro Sekunde durchprobieren. In der Praxis verwendet man deshalb absichtlich langsame Verfahren wie bcrypt oder Argon2 und ergänzt jedes Passwort vor dem Hashen mit einem zufälligen Wert (\textit{Salt}), damit gleiche Passwörter unterschiedliche Hashwerte erhalten. InstaHub speichert seine Passwörter zum Beispiel mit bcrypt.
    \begin{myanswer}
        \begin{lstlisting}
from sqlite3 import *
import hashlib

# ...

    password = input("Enter password")    # Ihr Klartext-Passwort zum "hashen"
    hash_object = hashlib.sha1()          # SHA-1-Hash-Objekt erstellen
    hash_object.update(password.encode()) # Hash-Objekt updaten mit der Byte-Darstellung des Passworts
    password = hash_object.hexdigest()    # Hexadezimale Darstellung erhalten

    # Restlicher Code
\end{lstlisting}
    \end{myanswer}
\end{mychallenge}

\begin{myexercise}[Users-Tabelle mit Game verknüpfen]
    Erweitern Sie Ihr Python-Game so, dass man zu Beginn aufgefordert wird, ein gültiges Login einzugeben. Erst wenn dies passiert ist, startet das Spiel. Sie können hierzu folgende Vorlage vervollständigen:
    \begin{lstlisting}
while True:
    username = input("Enter username")

    # TODO Den User mit diesem Benutzernamen in der Tabelle users suchen

    if True: # TODO Bedingung anpassen: Wenn der User gefunden wurde…
        break
    else:
        print("Unknown user. Try again.")

while True:
    password = input("Enter password")

    if True: # TODO Bedingung anpassen: Wenn das eingegebene Passwort stimmt…
        break
    else:
        print("Incorrect password. Try again.")
\end{lstlisting}
    \begin{myanswer}
        \begin{lstlisting}
from sqlite3 import *

with connect("game.db") as con:
    cursor = con.cursor()

    while True:
        username = input("Enter username")

        # Den User mit diesem Benutzernamen in der Tabelle users suchen
        sql = """
            SELECT *
            FROM users
            WHERE user = ?
        """
        cursor.execute(sql, (username,))
        resultset = cursor.fetchall()

        if len(resultset) > 0: # Wenn der User gefunden wurde...
            break
        else:
            print("Unknown user. Try again.")

    while True:
        password = input("Enter password")

        if resultset[0][2] == password: # Wenn das eingegebene Passwort stimmt
            break
        else:
            print("Incorrect password. Try again.")

    print("Welcome", username)

\end{lstlisting}
    \end{myanswer}
\end{myexercise}

\begin{myexercise}[\texttt{scores}-Tabelle mit Game verknüpfen]
    Schreiben Sie ein Programm zur Erstellung der \texttt{scores}-Tabelle und erweitern Sie Ihr Game so, dass am Ende jedes Spiels die Punktzahl darin abgespeichert wird.

    \textbf{Tipp}: Verwenden Sie den SQL-Befehl \lstinline[style=sql]|DATETIME('NOW')|, um die aktuelle Datums- und Zeit-Angabe zu erhalten.
    \begin{myanswer}
        \begin{lstlisting}
# ...

with connect("game.db") as con:
    # Datenbank-Cursor (=Lese/Schreibestift) holen
    cursor = con.cursor()

    # Punktzahl speichern
    sql = """
        INSERT INTO scores (user_id, score, date_score)
        VALUES
            (?, ?, DATETIME('NOW'))
    """
    cursor.execute(sql, (user_id, score))

\end{lstlisting}
    \end{myanswer}
\end{myexercise}

\begin{myexercise}[Highscore-Ansicht hinzufügen]
    Erweitern Sie Ihr Game mit einer Highscore-Ansicht, in welcher die zehn besten Punktestände inklusive Benutzernamen und Datum angezeigt werden.
    \begin{myanswer}
        \begin{lstlisting}
from sqlite3 import *

with connect("game.db") as con:
    cursor = con.cursor()

    #TODO TOP 10 scores inkl. Userdaten auslesen
    sql = """
        SELECT *
        FROM scores s
        JOIN users u ON u.id = s.user_id
        ORDER BY s.score DESC
        LIMIT 10
    """

    cursor.execute(sql)
    resultset = cursor.fetchall()

    for i in range(len(resultset)):
        hs = resultset[i]
        print(i+1, " User:", hs[5], "    Points:", hs[2])
        #TODO Mit PyGame-Befehl pg.draw.text(...) den Highscore auf dem Bildschirm anzeigen
\end{lstlisting}
    \end{myanswer}

\end{myexercise}

\begin{mychallenge}
    Erstellen Sie einen Super-User und verändern Sie Ihr Spiel so, dass die Grafik verändert wird, wenn der Super-User eingeloggt ist (z.B. wird ein goldener Ball statt einem normalen Ball angezeigt). Lassen Sie Ihrer Kreativität freien Lauf!
\end{mychallenge}

\begin{mychallenge}
    Zeigen Sie bei jedem Game-Start einen Text mit dem bisher höchsten Punktestand an, um den User dazu zu animieren, den Highscore zu knacken!
\end{mychallenge}